\chapter{Input Provider Framework}
\label{chap:input-providers}

The input provider framework governs how gert resolves dynamic input values
at run time.  When a runbook declares an input with a \texttt{from:} binding,
the provider framework locates the appropriate provider, dispatches a
resolution request, and returns the concrete value before the first step that
needs it begins.

\section{Design Goals}

\begin{itemize}
  \item \textbf{Extensibility}: third-party systems (PagerDuty, ServiceNow,
        Vault, etc.) integrate as provider plugins without changes to the host
        binary.
  \item \textbf{Transparency}: every resolved value is recorded in the run
        trace so the operator can audit where each input came from.
  \item \textbf{Fallback safety}: resolution failures degrade gracefully to an
        interactive prompt rather than aborting the run silently.
  \item \textbf{Consistency with tools}: providers share the same transport
        layer as tools (\texttt{stdio}, \texttt{stdio-jsonrpc}, \texttt{mcp}),
        minimising implementation surface.
\end{itemize}

\section{Provider Definition Schema}
\label{sec:provider-definition-schema}

Each provider is declared in a \texttt{.provider.yaml} file.  By convention,
project-level providers live in \texttt{providers/<name>.provider.yaml}.

\begin{verbatim}
# providers/pagerduty.provider.yaml
apiVersion: provider/v2

meta:
  name: pd                         # Prefix used in from: bindings (e.g. pd.*)
  version: "1.0.0"
  description: "Resolves inputs from PagerDuty incidents and services"
  binary: pd-provider              # Executable name (looked up in PATH)

transport:
  mode: stdio-jsonrpc              # stdio | stdio-jsonrpc | mcp
  startup:
    argv: ["--config", "pd.yaml"]
    ready-pattern: "provider ready"
    timeout: "10s"
    shutdown-method: "shutdown"

resolution:
  method: provider/resolve         # JSON-RPC method the host calls
  prefixes:
    - "pd."                        # All from: values starting with pd.
                                   # are dispatched to this provider

  # Schema for what the provider needs in each resolution request
  input-schema:
    type: object
    properties:
      bindings:
        type: object
        description: "Map of input name to from: binding string"
      context:
        type: object
        description: "Execution context (runId, userId, mode)"
    required: ["bindings"]

  # Schema for what the provider returns
  output-schema:
    type: object
    properties:
      resolved:
        type: object
        description: "Map of input name to resolved string value"
      warnings:
        type: array
        items:
          type: string
    required: ["resolved"]

  cache:
    enabled: true
    scope: run                     # run | step | none
    ttl: "60s"

governance:
  requires-capabilities:
    - capability/network           # This provider needs outbound network
  env-read-patterns:
    - "PD_API_KEY"
    - "PD_*"
\end{verbatim}

\section{Resolution Protocol}
\label{sec:provider-protocol}

\subsection{Resolution Flow}

When the run engine encounters an input with a \texttt{from:} binding (e.g.,
\texttt{from: pd.incident.title}), the following steps occur:

\begin{enumerate}
  \item \textbf{Prefix match}: the framework scans registered providers for one
        whose declared prefix matches the start of the binding string.  The
        longest matching prefix wins.
  \item \textbf{Provider start}: if the matched provider's process is not yet
        running, the host spawns it and waits for the ready signal (same
        mechanics as \texttt{stdio-jsonrpc} tools; see
        Section~\ref{sec:tool-transport}).
  \item \textbf{Batch collection}: all inputs whose binding prefix maps to the
        same provider are grouped into a single resolution request to minimise
        round-trips.
  \item \textbf{JSON-RPC call}: the host sends a \texttt{provider/resolve}
        request.
  \item \textbf{Result merge}: resolved values are merged into the run's
        variable map.
  \item \textbf{Fallback}: any binding that was not resolved (absent from the
        response's \texttt{resolved} map, or where the provider returned a
        warning) falls back according to the input's \texttt{fallback:} field
        (default: interactive prompt).
\end{enumerate}

\subsection{Resolution Request Envelope}

\begin{verbatim}
{
  "jsonrpc": "2.0",
  "id": "prov-1",
  "method": "provider/resolve",
  "params": {
    "bindings": {
      "incident_title": "pd.incident.title",
      "incident_severity": "pd.incident.severity",
      "service_name": "pd.service.name"
    },
    "context": {
      "runId":   "run-abc123",
      "stepId":  "__pre-flight__",
      "userId":  "ops@acme.example",
      "mode":    "real",
      "traceId": "trace-xyz"
    }
  }
}
\end{verbatim}

\subsection{Resolution Response Envelope}

\begin{verbatim}
{
  "jsonrpc": "2.0",
  "id": "prov-1",
  "result": {
    "resolved": {
      "incident_title": "Database connection pool exhausted",
      "incident_severity": "P1",
      "service_name": "payments-api"
    },
    "warnings": [],
    "telemetry": {
      "startedAt":  "2026-04-18T10:00:00Z",
      "finishedAt": "2026-04-18T10:00:00.080Z",
      "durationMs": 80
    }
  }
}
\end{verbatim}

\subsection{Resolution Error Handling}

If the provider returns a JSON-RPC error, or if a binding key is absent from
the \texttt{resolved} map, the host applies the input's fallback strategy:

\begin{verbatim}
# In runbook meta.inputs
inputs:
  incident_title:
    from: pd.incident.title
    description: "Active incident title"
    fallback: prompt              # prompt | fail | default
    default: "Unknown incident"   # Used when fallback: default
\end{verbatim}

\begin{itemize}
  \item \texttt{fallback: prompt} (default): the host issues an interactive
        prompt for the value.  In non-interactive modes (dry-run, CI) this
        causes the run to fail.
  \item \texttt{fallback: default}: the declared \texttt{default:} value is
        used without prompting.  A warning is written to the trace.
  \item \texttt{fallback: fail}: the run aborts immediately with a structured
        error.
\end{itemize}

\subsection{Resolution Caching}

Providers may declare a cache policy in their definition.  The host respects
this policy:

\begin{itemize}
  \item \texttt{scope: run} — the resolved value is reused for all subsequent
        steps in the same run that reference the same binding.  This is the
        default for external providers.
  \item \texttt{scope: step} — the value is resolved fresh on each step that
        references it.
  \item \texttt{scope: none} — caching is disabled; every binding triggers a
        new RPC call.
\end{itemize}

Cache entries are keyed by \texttt{(runId, bindingString)}.  Cached values are
stored in the run's in-memory variable map only; they are not persisted across
run resumptions.

\section{Built-in Providers}
\label{sec:builtin-providers}

The following providers are compiled into the gert host binary and are always
available without a \texttt{.provider.yaml} file.

\subsection{env}

Reads values from environment variables.

\begin{verbatim}
inputs:
  api_key:
    from: env.ACME_API_KEY
    description: "API key for the Acme service"
\end{verbatim}

The \texttt{env.} prefix is followed by the exact environment variable name.
If the variable is not set, resolution fails and the fallback strategy is
applied.  The \texttt{env} provider never contacts an external process; it
reads directly from the host's environment.

\subsection{file}

Reads values from files on the local file system.

\begin{verbatim}
inputs:
  ssh_key:
    from: "file.~/.ssh/id_rsa"
    description: "SSH private key"
  config_value:
    from: "file./etc/gert/config.json#/database/host"
    description: "Database host from config file"
\end{verbatim}

The path after \texttt{file.} is the file path.  An optional JSON Pointer
fragment (\texttt{\#/path/to/field}) extracts a nested value from a JSON or
YAML file.  Without a fragment, the entire file contents are returned as a
string.

\subsection{prompt}

The interactive prompt provider is the default resolution strategy.  It is
also the fallback for any failed provider resolution (unless overridden).

\begin{verbatim}
inputs:
  server_name:
    from: prompt
    description: "Server to check"
\end{verbatim}

In \texttt{real} and \texttt{debug} modes, the host presents a terminal prompt
to the operator.  In \texttt{dry-run} and \texttt{replay} modes, the prompt
provider reads from the scenario's \texttt{inputs.yaml} instead.  In
non-interactive environments where no scenario is available, the prompt provider
fails with an informative error directing the operator to supply the value via
\texttt{--var} or an explicit provider.

\subsection{workspace}

Reads values from the workspace configuration file
(\texttt{.gert/config.yaml}).

\begin{verbatim}
inputs:
  default_region:
    from: workspace.defaults.region
    description: "AWS region from workspace config"
\end{verbatim}

The path after \texttt{workspace.} is a dot-notation key into the workspace
configuration object.  This allows teams to centralise defaults without
repeating them in every runbook.

\section{Interactive Step Type Contracts}
\label{sec:interactive-step-contracts}

Input providers not only resolve \texttt{from:} bindings for runbook inputs,
they also power the three interactive step types: \texttt{choice},
\texttt{decision}, and \texttt{collector}.  When the run engine executes one of
these steps, it dispatches a step-specific request to the configured input
provider (or falls back to the built-in \texttt{prompt} provider).

\subsection{Choice Step Contract}
\label{subsec:choice-contract}

A \texttt{choice} step presents the operator with a fixed set of labeled
options and stores the selected value as a variable.  This is analogous to a
dropdown menu or radio button group in a graphical form.

\paragraph{Request envelope.}
When the run engine executes a \texttt{choice} step, it sends a
\texttt{provider/choice} JSON-RPC request to the configured provider:

\begin{verbatim}
{
  "jsonrpc": "2.0",
  "id": "choice-1",
  "method": "provider/choice",
  "params": {
    "stepId": "select-region",
    "prompt": "Select the target AWS region",
    "options": [
      { "value": "us-east-1", "label": "US East (N. Virginia)" },
      { "value": "us-west-2", "label": "US West (Oregon)" },
      { "value": "eu-west-1", "label": "Europe (Ireland)" }
    ],
    "default": "us-east-1",
    "context": {
      "runId": "run-abc123",
      "stepId": "select-region",
      "userId": "ops@acme.example",
      "mode": "real"
    }
  }
}
\end{verbatim}

\paragraph{Response envelope.}
The provider returns the selected value:

\begin{verbatim}
{
  "jsonrpc": "2.0",
  "id": "choice-1",
  "result": {
    "selected": "us-west-2"
  }
}
\end{verbatim}

\paragraph{Contract requirements.}
\begin{itemize}
  \item The \texttt{options} array must contain at least two elements.
  \item Each option must have both \texttt{value} (the string stored in the
        variable) and \texttt{label} (the human-readable text displayed to the
        operator).
  \item The \texttt{default} field is optional.  If present, it must match one
        of the \texttt{value} fields in the \texttt{options} array.
  \item The provider's \texttt{selected} response must match one of the option
        values exactly.  If the provider returns a value not in the options
        list, the run engine fails the step with a
        \texttt{choice.invalid\_selection} error.
  \item Providers MAY render options in the order provided, or MAY sort them
        alphabetically by label for improved usability.  The canonical order is
        declared in the runbook YAML.
\end{itemize}

\paragraph{Built-in prompt provider behaviour.}
The built-in \texttt{prompt} provider presents a numbered list in the terminal
and waits for the operator to enter the number or value corresponding to their
selection.  In \texttt{dry-run} mode, it reads the value from the scenario's
\texttt{inputs.yaml} or uses the \texttt{default} if provided.

\subsection{Decision (Router) Step Contract}
\label{subsec:decision-contract}

A \texttt{decision} step (also known as \texttt{router} in v1) presents the
operator with a set of routing choices, where each choice directs execution to
a different runbook or labeled section.  Unlike \texttt{choice}, a
\texttt{decision} is a control-flow primitive: the selected route determines
which node the run engine transitions to next.

\paragraph{Request envelope.}
When the run engine executes a \texttt{decision} step, it sends a
\texttt{provider/decision} JSON-RPC request:

\begin{verbatim}
{
  "jsonrpc": "2.0",
  "id": "decision-1",
  "method": "provider/decision",
  "params": {
    "stepId": "triage-severity",
    "prompt": "Select the incident severity level",
    "routes": [
      { "route": "p1-critical", "label": "P1: Critical (immediate)" },
      { "route": "p2-high",     "label": "P2: High (1 hour SLA)" },
      { "route": "p3-normal",   "label": "P3: Normal (24 hour SLA)" }
    ],
    "context": {
      "runId": "run-abc123",
      "stepId": "triage-severity",
      "userId": "ops@acme.example",
      "mode": "real"
    }
  }
}
\end{verbatim}

\paragraph{Response envelope.}
The provider returns the selected route label:

\begin{verbatim}
{
  "jsonrpc": "2.0",
  "id": "decision-1",
  "result": {
    "route": "p2-high"
  }
}
\end{verbatim}

\paragraph{Contract requirements.}
\begin{itemize}
  \item The \texttt{routes} array must contain at least two elements.
  \item Each route must have both \texttt{route} (the target node label or
        runbook ID) and \texttt{label} (the human-readable text).
  \item The provider does \emph{not} need to know about the runbook graph
        structure.  It only returns the selected route label.  The run engine
        is responsible for resolving the route label to a target node and
        validating that the target exists.
  \item If the provider returns a route label not in the \texttt{routes} array,
        the run engine fails the step with a \texttt{decision.invalid\_route}
        error.
  \item The decision result is \emph{not} stored as a variable.  It only
        controls execution flow.  If the operator needs the decision recorded,
        the runbook author must add an explicit variable assignment in the
        target node.
\end{itemize}

\paragraph{Built-in prompt provider behaviour.}
The built-in \texttt{prompt} provider displays the routes as a numbered list
and waits for the operator to select one.  In \texttt{dry-run} mode, the route
is read from the scenario file or the run fails if no scenario is provided.

\subsection{Collector Step Contract}
\label{subsec:collector-contract}

A \texttt{collector} step gathers unstructured input from the operator: free-form
text, multi-line notes, file attachments, screenshots, URLs, or other evidence
artifacts.  The collected data is stored as named variables or artifacts in the
run's state.

\paragraph{Request envelope.}
When the run engine executes a \texttt{collector} step, it sends a
\texttt{provider/collect} JSON-RPC request:

\begin{verbatim}
{
  "jsonrpc": "2.0",
  "id": "collect-1",
  "method": "provider/collect",
  "params": {
    "stepId": "gather-evidence",
    "instructions": "Upload a screenshot of the error and describe the issue",
    "fields": [
      {
        "name": "screenshot",
        "type": "file",
        "label": "Error screenshot",
        "required": true,
        "accept": ["image/png", "image/jpeg"],
        "maxSizeBytes": 10485760
      },
      {
        "name": "description",
        "type": "text",
        "label": "Describe what you observed",
        "required": true,
        "multiline": true
      },
      {
        "name": "log_url",
        "type": "url",
        "label": "Link to related logs (optional)",
        "required": false
      }
    ],
    "context": {
      "runId": "run-abc123",
      "stepId": "gather-evidence",
      "userId": "ops@acme.example",
      "mode": "real"
    }
  }
}
\end{verbatim}

\paragraph{Response envelope.}
The provider returns the collected values and any uploaded artifacts:

\begin{verbatim}
{
  "jsonrpc": "2.0",
  "id": "collect-1",
  "result": {
    "values": {
      "description": "Database connection pool exhausted...",
      "log_url": "https://logs.acme.example/trace/xyz"
    },
    "artifacts": [
      {
        "field": "screenshot",
        "filename": "error-2026-04-18.png",
        "contentType": "image/png",
        "sizeBytes": 245631,
        "sha256": "a3c7d1f...",
        "storagePath": "artifacts/run-abc123/screenshot.png"
      }
    ]
  }
}
\end{verbatim}

\paragraph{Contract requirements.}
\begin{itemize}
  \item The \texttt{fields} array declares the form structure.  Each field has:
        \begin{itemize}
          \item \texttt{name}: the variable name where the value will be stored.
          \item \texttt{type}: one of \texttt{text}, \texttt{file}, \texttt{url},
                \texttt{number}, or \texttt{date}.
          \item \texttt{label}: human-readable prompt text.
          \item \texttt{required}: boolean flag.
          \item Type-specific constraints: \texttt{accept} (MIME types for file
                uploads), \texttt{maxSizeBytes} (file size limit), \texttt{multiline}
                (for text fields).
        \end{itemize}
  \item Providers MUST return all required fields in the \texttt{values} map.
        If a required field is missing, the run engine fails the step with a
        \texttt{collector.missing\_required\_field} error.
  \item For \texttt{type: file} fields, providers MUST upload the file to the
        artifact store and return metadata in the \texttt{artifacts} array.
        The artifact metadata includes:
        \begin{itemize}
          \item \texttt{field}: the field name from the request.
          \item \texttt{filename}: the original filename.
          \item \texttt{contentType}: MIME type.
          \item \texttt{sizeBytes}: file size in bytes.
          \item \texttt{sha256}: SHA-256 hash for integrity verification.
          \item \texttt{storagePath}: path where the artifact is stored (relative
                to the run's artifact root).
        \end{itemize}
  \item File size limits are enforced by the provider.  If an uploaded file
        exceeds \texttt{maxSizeBytes}, the provider MUST reject it with a
        \texttt{collector.file\_too\_large} error before storing it.
  \item Supported MIME types for file uploads:
        \begin{itemize}
          \item Images: \texttt{image/png}, \texttt{image/jpeg}, \texttt{image/gif}
          \item Documents: \texttt{application/pdf}, \texttt{text/plain},
                \texttt{text/markdown}
          \item Archives: \texttt{application/zip}, \texttt{application/x-tar},
                \texttt{application/gzip}
          \item Logs: \texttt{text/plain}, \texttt{application/json}
        \end{itemize}
  \item Providers MAY support additional MIME types, but MUST respect the
        \texttt{accept} constraint if declared in the field definition.
  \item The artifact storage path is provider-specific.  The built-in
        \texttt{prompt} provider stores artifacts in
        \texttt{.gert/runs/<runId>/artifacts/}.  External providers (e.g., S3,
        Azure Blob) MAY store artifacts in remote storage and return a signed
        URL or storage reference.
\end{itemize}

\paragraph{Built-in prompt provider behaviour.}
The built-in \texttt{prompt} provider displays each field as a terminal prompt.
For \texttt{type: text} with \texttt{multiline: true}, it opens the operator's
configured editor (\texttt{\$EDITOR}).  For \texttt{type: file}, it prompts for
a file path and copies the file to the local artifact store.  In
\texttt{dry-run} mode, all values are read from the scenario file.

\subsection{Provider Capability Matrix}
\label{subsec:provider-capability-matrix}

Not all providers support all interactive step types.  The following table
shows which step types each built-in provider supports:

\begin{center}
\begin{tabular}{|l|c|c|c|c|}
\hline
\textbf{Provider} & \textbf{Input Resolution} & \textbf{choice} & \textbf{decision} & \textbf{collector} \\
\hline
\texttt{prompt}   & Yes & Yes & Yes & Yes \\
\texttt{env}      & Yes & No  & No  & No  \\
\texttt{file}     & Yes & No  & No  & No  \\
\texttt{workspace}& Yes & No  & No  & No  \\
\hline
\end{tabular}
\end{center}

External providers (PagerDuty, ServiceNow, Slack, etc.) MAY implement any
combination of these contracts.  A provider that implements
\texttt{provider/choice} MUST declare this in its capability advertisement
during the handshake phase:

\begin{verbatim}
# Example provider capability declaration
{
  "jsonrpc": "2.0",
  "method": "provider/initialize",
  "params": {
    "capabilities": {
      "resolution": true,
      "choice": true,
      "decision": false,
      "collector": true
    }
  }
}
\end{verbatim}

If the run engine encounters a \texttt{choice}, \texttt{decision}, or
\texttt{collector} step and the configured provider does not advertise the
corresponding capability, the engine falls back to the built-in \texttt{prompt}
provider.

\section{Provider Composition}
\label{sec:provider-composition}

A \texttt{from:} binding may declare a priority-ordered chain of providers.
The host attempts each provider in order and uses the first successful result.

\begin{verbatim}
inputs:
  api_token:
    from:
      - env.ACME_API_TOKEN          # Try env var first
      - vault.secret/acme/api-token # Fall back to Vault
      - prompt                      # Final fallback: ask the operator
    description: "API token for Acme service"
\end{verbatim}

Chaining is evaluated left-to-right.  A provider is considered to have failed
if it returns an error, if the binding key is absent from the resolved map, or
if the resolved value is an empty string (unless the input declares
\texttt{allow-empty: true}).

\section{Provider Lifecycle}
\label{sec:provider-lifecycle}

\paragraph{Process model.}
Providers that use the \texttt{stdio-jsonrpc} or \texttt{mcp} transport are
\emph{persistent}: the host spawns each provider process once at the start of
the pre-flight resolution phase and keeps it alive for the entire run.  This
avoids repeated startup latency for providers that need to authenticate or
maintain session state (e.g., Vault token renewal, PagerDuty OAuth).

Providers using the \texttt{stdio} transport are \emph{per-resolution}: the
host spawns a new process for each resolution request.

\paragraph{Startup.}
Provider processes are started lazily on first use, following the same
\texttt{ready-pattern} / timeout mechanism as \texttt{stdio-jsonrpc} tools
(Section~\ref{sec:tool-transport}).

\paragraph{Shutdown.}
At run completion (success, failure, or cancellation) the host iterates over
all running provider processes and sends the method declared in
\texttt{transport.startup.shutdown-method} (default: \texttt{shutdown}).
Providers are given a 5-second grace period to clean up before the host sends
SIGTERM, then SIGKILL after a further 5 seconds.

\paragraph{Crash handling.}
If a provider process crashes during a resolution call, the host logs a
structured \texttt{provider.crashed} trace event and applies the input's
\texttt{fallback:} strategy for all bindings that were in-flight at the time
of the crash.
